\begin{chapterbox}
    \chapter{Die Legende der entzweiten Zeit (2023)}
    \label{Die Legende der entzweiten Zeit (2023)}
    \az{63 und 64}

    \begin{center}
        Fortsetzung von \hypref{Der Giftzwerg und die Sphäre (2020)}
    \end{center}
    
    Instabile Zeitportale verbinden zwei Zeiten: Das herbstliche Andor, dessen Bewohner noch nichts von der ewigen Kälte ahnen, und das winterliche Andor, wo einige letzte Helden die Rietburg tapfer verteidigen.
\end{chapterbox}


\az{Jahr 64}

Es war ein eiskalter Tag in Andor. Bauern bibberten in der Rietburg, Kreaturen schlichen durch den Schnee und die Vorratskammern der Andori wurden leerer und leerer. \bigskip

Die Ewige Kälte lag über Andor. Alle Helden waren ins Land der Steppe gezogen, um rettende Kräuter für die Eisschlafenden zu sammeln. Alle? Nein! Einige wenige Helden waren in Andor verblieben und wehrten eine Vielzahl von Kreaturen ab. So viele von ihnen auch dem Eisschlaf anfielen, stetig rückten neue Dunkle Gegner nach. Die Verteidiger Andors wurden immer weniger, ohne dass gute Neuigkeiten aus dem Osten kämen. 

Und noch wussten die Helden nicht, dass eine finstere Macht weitere Gefahren nach Andor rief. Gefahren, die sie mit Dunkler Magie in einem festen Bann hielt. \bigskip

Fürst Hallgard hatte eine Notkonferenz einberufen. Der Herrscher über Cavern spürte bereits die Kälte in seinen Knochen aufsteigen und befürchtete, dass es zu spät sei. „Unsere klügsten Köpfe können dem Winterstein mit ihrem komplexesten Konstruktionen keinen Kratzer zufügen, ohne zu Eisstatuen zu erstarren! Keiner von Rekas Tränken vermag die Kälte zu lindern! Ist dies unser Ende?“ „Warum schicken wir den Winterstein nicht woandershin, weit weg?“, fragte der grimmige Radan. 

„Keine Lehre aus Thoralds Taten gezogen?“, lachte Hallgard trocken. „Wir riskieren damit ohnehin nur, dass dieses mächtige Artefakt in feindliche Finger fällt.“

„Es fehlt uns an der Zeit“, murmelte Mart, „Wir würden bestimmt eine Lösung finden, wenn wir den Effekt des Artefakts nur für eine gewisse Zeit eindämmen könnten.“ „Das ist es! ZEIT!“, rief Hallgard und riss seine Augen auf. Der Zwergenfürst mahlte mit seinem Unterkiefer und schien mit sich selbst zu hadern. Dann rief er bestimmt: „Alle raus aus diesem Raum! Alle ... außer du, Mart!“ \bigskip

„Mart, höre mir gut zu“, flüsterte Hallgard, sobald alle anderen den Raum verlassen hatten. „Ich werde bald der Ewigen Kälte zum Opfer fallen. Meine Beine sind bereits taub geworden. Doch deine sind es nicht. Nimm diesen Schlüssel und eile zur geheimen Fürstenkammer. Die dort gesammelten Artefakte offenbaren Antworten auf das Leben, das Universum und den ganzen Rest. Ein bestimmtes Artefakt kann uns hoffentlich in der hiesigen Lage weiterhelfen. Reise zur Geheimen Kammer, öffne sie mit meinem Fürstenschlüssel und bringe mir die Sphäre zurück. Eine faustgroße metallene Kugel aus einem Geflecht ineinander verschlungener Streben. Du wirst sie erkennen, wenn du sie siehst.“ 

„Und was soll ich mit dieser Sphäre tun?“, fragte Mart. „Du? Nichts!“, keuchte Hallgard, „Das Wissen um diese Apparatur ist zu gefährlich, als dass es irgendjemandem mehr anvertraut werden sollte, als zwingend nötig. Ich selbst werde die Sphäre handhaben. Bringe sie mir, ehe auch ich dem Eisschlaf anheimfalle!“ \bigskip

Während Mart mit dem Schlüssel zur geheimen Fürstenkammer davonhastete, ließ Hallgard sich auf seinen Thron zurücksinken. Seine Gedanken rasten. Hätte er Mart doch vom Geheimnis der Sphäre erzählten sollen? Hallgard spürte, wie die Ewige Kälte durch sein Fleisch kroch und seine Muskeln verhärtete.

Wohl ahnend, dass er nicht mehr atmen würde, wenn Mart zurückkehrte, stolperte der Fürst zu seinem Schreibpult. Hektisch tauchte er eine silberne Feder in ein Fässchen Geheimtinte, welche in wenigen Tagen verblassen würde. Dann suchte er in seinem Gedächtnis die Instruktionen hervor, die ihm sein Vater Hallwort einst hinterlassen hatte, ehe dieser neugierig in den Norden gereist war. Seinem Tod entgegen.

Hallgard wünschte sich, er hätte damals besser zugehört. Doch so viele andere Gedanken waren zu dieser Zeit durch seinen noch nicht einmal 80 Jahre jungen Kopf geschossen, allen voran Wut auf seinen Vater und dessen Missachtung der Zwergentradition.

Was hatte Hallwort schon wieder gesagt, wie man das Ding aktivierte? Und worauf sollte man unbedingt achten, wenn man es einsetzte? \bigskip

Mart hatte die Hoffnung nie aufgegeben. Die schwach summende silberne Sphäre in seiner Manteltasche versprach Rettung für Cavern und die umliegenden Reiche. Doch die Ansammlung trauernder Zwerge vor dem Thronsaal verriet Mart, dass er nicht schnell genug gewesen war. Auch den weisen Fürsten hatte es erwischt. Man ließ Mart nicht einmal den erstarrten Hallgard sehen. Mit gesenktem Kopf reichte ihm seine Schwester Ragga, eine Wächterin Hallgards, ein versiegeltes Pergament. „Vom Fürsten. Nur für deine Augen bestimmt“, sprach sie. Mit zitternden Fingern entrollte Mart die Botschaft. Zittrige Runen und hastige Skizzen übersäten sie. Instruktionen von Hallgard. In schlichten Lettern stand dort beschrieben, wie man die Streben und Schalter zu betätigen hatte, um die Sphäre zu nutzen und – nein, Hallgard schien Mart immer noch nicht verraten zu wollen, was die Apparatur denn tat, wenn sie aktiviert worden war. Doch hatte er ihm hoffentlich genug mitgeteilt, dass Mart es selbst herausfinden konnte. 

„Das Schicksal liegt in deinen Händen, Mart“, endete die Nachricht des Fürsten. „Ich vertraue auf deine Diskretion. Denn, wenn diese Macht in die falschen Hände gerät, ist noch so viel mehr als nur das jetzige Cavern in Gefahr.“ \bigskip

Ohne zu verstehen, kehrte Mart rasch in seine Wohnkammer zurück, an seinen zwei Eulenküken vorbei, in seinen kleinen Schlafraum. Dort versuchte er im flackernden Licht eines Glutholzes, Hallgards zu folgen. Doch waren sie unvollständig, widersprüchlich, unverständlich. Dazu kam, dass die Zeichen auf dem Dokument stetig verblassten. Wollte er sie abschreiben? Durfte er dies überhaupt?

Mart schob die letzte Querverstrebung runter und legte ein Kontrollzelle in der Mitte der Kugel frei. Doch konnte er nicht mehr entziffern, welchen Hebel er nun umlegen musste. Wagte er es, zu improvisieren? \bigskip

Da klopfte es an der Tür. Herein trat Marts treuer Lebensgefährte Bort, mit einer Tasse warmen Tees in der Hand und bestimmt einigen beruhigenden Worten auf der Zungenspitze. Doch statt diese Worte zu teilen, erstarrte Bort bloß, als er die Sphäre in Marts Hand erblickte.
„Bei Nehals feurigem Kot! Woher hast du die?“

„Aus der geheimen Kammer unseres Fürsten. Weißt du etwa, was sie ist? Was sie kann?“

„Ein wenig. Ich bin ihr bloß einmal begegnet, und das ist lange her.“ Bort starrte die Sphäre ängstlich an, ehe er anhängte: „Es verwundert mich, dass sie ihren Weg zurück nach Cavern fand. Sei vorsichtig mit ihr, Mart.“

„Bin ich doch stets“, versuchte Mart, selbstsicher zu erscheinen.

„Kaulquappenquatsch!“, entgegnete Bort. 

„Hallwort nannte die Sphäre unsere letzte Hoffnung.“ „Dem mag so sein“, murmelte Bort, „Doch habe ich keine guten Erinnerungen an sie. Wir mussten beim feurigen Gott schwören, niemandem von der Wirkungsweise dieser Apparatur zu erzählen, damit die Details nicht in die falschen Hände gelangen konnten.“

„Du musst mir ja nichts erzählen, nur helfen. Erinnerst du dich zufälligerweise noch daran, wie die Sphäre ausgelöst wird? Die letzten Aufzeichnungen unseres Fürsten sind leider bloß lückenhaft – und werden immer blasser.“

Bort kratzte sich am Bart. „Kaum. Ich war noch ein frischgebackener Wächter und löste die Sphäre nie selbst aus. Darf ich?“

Vorsichtig überreichte Mart seinem Partner die verblassenden Pläne. Bort verglich die Verstrebungen der Sphäre mit den Skizzen und bemängelte einige Unterschiede, welche auch Mart schon aufgefallen waren. „Nein, das kommt mir alles unbekannt vor. Warte, ich hole Papier“, meinte Bort, „Besser, wir zeichnen ab, was wir an Anweisungen haben, ehe sie völlig verschwinden.“

Schon war er wieder zum Zimmer hinausgehuscht. Mart blickte angestrengt auf Hallgards Notizen und versuchte sie sich einzuprägen. \bigskip

„Kreatok, steh‘ mir bei!“, erklang Borts fluchende Stimme aus dem nächsten Raum, gefolgt von einem lauten Scheppern. Vermutlich war Bort mal wieder zu hastig gerannt und hatte einen Tisch gestreift. 

„Alles grün?“, erkundigte sich Mart. Keine Antwort ertönte. Das war leicht besorgniserregend.

„Was machst du da?“, erklang eine heisere Stimme direkt hinter Mart, die nicht zu Bort gehörte. Mart fuhr herum und erstarrte. Das war definitiv besorgniserregend. Vor ihm schwebte die bläulich schimmernde Gestalt eines Greises, der griesgrämig auf ihn herunterblickte. Ein Mensch, allem Anschein nach. Oder zumindest war er einst einer gewesen. Sein gebückter Kopf streifte den Türrahmen. Hinter ihm konnte Mart knapp eine am Boden liegende Gestalt ausmachen, die unschön wie Bort aussah. „Nein, nein, hör nicht auf! Das sieht ja höchst interessant aus, was du hier tust“, sprach der Greis. „Ich kann mich nur nicht entscheiden, ob ich dieses Artefakt aufhalten oder studieren sollte.“ 

„Wer ... wer seid Ihr?“, fragte Mart.

„Das ist kaum von Belang. Ich bin hier, um euch zu helfen“, nickte der Greis, „Und diese Sphäre sieht so aus, als könnte sie mir sehr gut helfen. Wenn sie nicht falsch ausgelöst wird. Ich vermute ...
“
Er starrte in die Kugel hinein.

„... ja, ich vermute, du musst nur noch dieses Zahnrad dahinten zwei Speichen weiterdrehen. Dann kommt alles gut.“ 

In Gegenwart dieses eigenartigen Wesens weiter an diesem streng geheimen und höchst gefährlichen Artefakt zu werkeln, erschien Mart alles andere als angebracht. Welche anderen Optionen standen ihm offen? 

„Dann mache ich es halt selbst“, grummelte der Greis. Er langte nach der Sphäre. Da hatte Mart aber genug, sprang selbst vor und schnappte dem Greis die Sphäre weg.


Ein wenig zu energisch. 

Ein lautes Klicken ertönte.

Mart hatte einen Mechanismus ausgelöst.

Die Sphäre begann zu surren und zu dampfen. Blauer Dampf stieß aus einer Öffnung in ihrer Seite, waberte in die Luft und wogte, festigte sich jedoch nicht. Ein Knall ertönte. Alles wurde dunkel. Mart fühlte, wie die Welt sich um ihn herum drehte, ihm die Luft ausging ... \bigskip

\az{Jahr 63}

... und frische Luft in seine Lungen strömte, während er nach Atem japste. Er blickte sich ängstlich um. Er lag auf dem Boden seiner Wohnhöhle, doch etwas war falsch. Der flauschige Teppich, auf dem Marts Gesicht ruhte, den sollte es eigentlich nicht mehr geben. Der war bei einem feurigen Unfall vor einigen Wochen verbrannt, und Mart hatte sich bislang nicht um einen Ersatz gekümmert. 

Vorsichtig erhob sich Mart und blickte sich in seiner Wohnhöhle um. Da war noch mehr, das nicht stimmte. Sein Schaukelstuhl und sein Ecktisch standen beide am falschen Ort, noch wie von vor seiner kürzlichen großen Höhlenumstellungsaktion. Der mächtige rote Edelstein, den Mart unlängst erworben hatte, hing nicht an der Wand. Bort war nirgends mehr zu sehen. Dafür war da immer noch dieser bläulich glühende Greis, der auf ihn herunterstarrte.

„Was hast du getan?!“, rief der Greis. Wütend grabschte er die Sphäre aus Marts Griff und starrte sie an. Das Ding stotterte und zuckte vor sich hin und reagierte nicht auf die Versuche des Greises, etwas an ihr einzustellen. Sein Mund verzog sich und er brüllte erneut: „Was hast du getan?! Die Schatten der Gegenwart und der Zukunft ... sie sind alle falsch! Verzerrt! Verändert!“

Mart wusste selbst nicht, was genau er getan hatte. Musste er auch nicht. Denn in diesem Augenblick beruhigte sich die Gestalt schlagartig. Breit grinsend murmelte sie: „Ah, natürlich! Ich machte mir umsonst Sorgen. Wie konnte ich bloß so blind sein? Wie konnte ich das nur vergessen?“ Ohne weitere Erklärungen breitete der Greis gebieterisch seine Arme aus und rief einige finstere Worte in einer fremden Sprache. Dunkelheit überkam Mart und er sank in einen tiefen, traumlosen Schlaf. \bigskip

Ja, Mart war wahrlich durch diese mysteriöse Sphäre in die Vergangenheit gereist, in einen Herbst, in dem noch niemand etwas von der Ewigen Kälte ahnte. Und die stotternde Sphäre sorgte nicht nur beim glühenden Greis für Sorgen. Schwirrende, instabile bläuliche Portale öffneten sich spontan an verschiedenen Stellen im Lande und verbanden die zwei Zeiten: Das herbstliche Andor, in das Mart und der Greis gereist waren, dessen Bewohner noch nichts von der ewigen Kälte ahnen, und das winterliche Andor, aus dem Mart und der Greis gereist waren, wo einige letzte Helden die Rietburg tapfer verteidigen. Nachdem Mart wieder erwacht war, suchte er zögerlich die herbstlichen Helden und berichtete ihnen: 

Ein gruseliger, bläulicher Greis habe ihn überwältigt und ein wertvolles Artefakt der Schildzwerge gestohlen! Es war nicht auszuschließen, dass dieser Vorfall für diese eigenartigen Zeitportale verantwortlich war, welche durch Andor waberten. Obwohl Mart sich geheimnisvoll gab und keine weiteren Details zu diesem gestohlenen Artefakt herausrückte, war klar: Die Helden wollten ihm helfen und den Greis aufspüren. 

\az{63 und 64}

So kam es, dass, während die meisten noch nicht eisschlafenden Helden im Land der drei Brüder den Dunklen Magier Orweyn suchten, von Orweyn prophezeite Bedrohungen überwanden und Wintersteine verwandelten – ein Abenteuer, das später als Legende der Pranken des Winters in die Chroniken der Bewahrer vom Baum der Lieder eingehen würde – auch in Andor nicht eisschlafende Helden den glühenden Greis jagten, und von ihm nach Andor geschickte Bedrohungen bannten. Dieses Abenteuer würde eines Tages als Legende der entzweiten Zeit in die Chroniken der Bewahrer vom Baum der Lieder eingehen. 

Mart würde später erfahren, wie sich Herbst-Helden aus der Vergangenheit mit Winter-Helden aus der Zukunft zusammentaten. Wie die Helden durch fünf in der Vergangenheit am Krallenfelsen versteckte und in der Zukunft am Krallenfelsen gefundene Goldstücke unmissverständlich klarstellten, dass diese Zeiten zusammenhingen. 

Diese Portale verbanden keine zwei Welten, sondern Vergangenheit und Zukunft derselben Zeitlinie, die sich aus irgendeinem Grund entzweit hatte. Zwei Zeiten, die sich parallel entwickelten, und von denen die eine dann doch die andere werden würde. Ein vorgeschriebenes Schicksal, scheinbar in Stein gemeißelt.

Für viel Verwirrung sorgte dies. Einzig die Kräuterhexe Reka hatte – dank einer lange zurückliegenden Beratung mit dem tulgorischen Hüter der Zeit – eine Ahnung, was es hiermit auf sich hatte. Die Helden erreichten den winterlichen glühenden Greis gerade rechtzeitig, um zu verhindern, dass er mit dem gestohlenen Winterstein von Cavern außer Landes fliehen konnte, und um zu verhindern, dass er die Nachricht über das Verwandeln des Wintersteins im Land der drei Brüder an sein früheres Selbst weitergeben konnte.\bigskip

Der Greis stürzte zu Boden. „Ich wusste, dass ihr mich aufhalten werdet. Eigenartig, diese Zeitgefüge. Wie kann es sein, dass ich mein früheres Selbst nicht warnen kann, obwohl ich weiß, was euch gleich erwartet?” Er hob seine Hände, blaue Stränge der Magie formend, bereit, erneut davonzuteleportieren ... und sackte bewusstlos zusammen, als ihn ein heldenhafter Schlag auf den Hinterkopf traf.\bigskip

Die Helden brachten die Gestalt zum Kerker unter der verschneiten Rietburg. Eine dortige Zelle hielt eine waschechte Schattenhexe gefangen, da würde doch hoffentlich auch ein teleportierender Greis nicht auszubrechen vermögen.\bigskip

Die Helden waren frohen Mutes, dass der ältere Tranuk und seine neuerdings schreckliche Steppenechse den Winterstein sicher ins Land der drei Brüder zurückbringen würden. Nun war es an der Zeit zu feiern. \bigskip

Ehe einige von ihnen in ihre eigene Zeit zurückkehren würden, versammelten sich alle beteiligten Helden, der jüngere Tranuk und einige andere Eingeweihte ein letztes Mal in der Vergangenheit, bei der herbstlichen Taverne zum Trunkenen Troll. Sie stießen an auf ihren Sieg über die Ewige Kälte, an einem Tisch ein wenig abseits des Hauses, um kein Misstrauen vergangener Tavernengäste zu wecken. Die beiden Rekas unterhielten sich prächtig und ließen sich zu einer Honigmilch überreden, auch wenn die ältere seltsam besorgt aussah. 

Ein Held bekam feuchte Augen, als er die jüngere Gilda in der Ferne eine ihrer Balladen schmettern sah. Bald würde die ältere aus dem Eisschlaf erwachen und wieder solche Lieder zum Besten geben können. 

Manch einem schmerzte der Kopf darüber, was sie soeben erlebt hatten. Oder war es Gildas Met, der ihnen zu Kopfe stieg? Hmmm ... schmeckte dieser Met nicht irgendwie bitterer als sonst?

Einer nach dem anderen kippten die Trinkenden um, als sie ihre Kräfte verließen. Manche rutschten gar von ihren Stühlen. 

Mit einem Knall erschien die jüngere glühende Gestalt triumphierend auf dem Tisch. „Das habt ihr davon, mich aus meinem Versteck zu verjagen versuchen! Ob ihr wohl ahntet, was ich vorhatte? Doch fürchtet euch nicht, ich werde euch nicht schaden. Ich bringe die Zeitverirrten unter euch sogar in ihre eigene Zeit zurück. Ihr Helden müsst es schließlich noch mit Varkur aufnehmen.“ „Wasss?“, nuschelte einer der Helden. 

Der jüngere Greis schnaubte. „Ich habe es vorausgesehen. Ihr werdet Qurun überwinden, das drohende Ende Hadrias, Ihr werdet die zankenden Orden vereinen. Ihr tut Gutes. Und so sollt ihr weiter leben und fröhlich walten. Dennoch will ich euch natürlich nicht die Gelegenheit geben, mich aufzuhalten. Darum dürft ihr euch nicht daran erinnern, was in den letzten Tagen geschah.“ 

„Wiiie?“, murmelte der Held.

„So!“, grinste der Greis. Er zückte eine Phiole, in der nur noch ein kleiner Rest eines leuchtenden Trunks umherschwappte. „Eine Tinktur gegen Seelenschmerz. In geringen Dosen kann sie schmerzhafte Erinnerungen erträglich machen oder gar ganz löschen. Doch vermag sie vieles mehr. Wenn ich sie richtig dosiert habe, werdet ihr gerade schön die Ereignisse der letzten Tage vergessen. Damit ich beruhigt das Eintreffen der Ewigen Kälte vorbereiten kann. Welches, wie ich weiß, ein gutes Ende für mich nehmen wird. Und ihr ... ihr werdet wohl in die sichere Ferne fliehen? Ich empfehle Hadria.“

Ein diabolisches Gekicher entfloh dem Greis.

Dann umfing die Helden der Mantel der Ungewissheit. \bigskip



\az{Jahr 64}

Auf dem Weg ins Land der drei Brüder empfing Tranuk im Traum eine seltsam klare Vision von einem geheimnisvollen Eiswesen, welches den Winterstein hungrig beobachtete. Am nächsten Tag fühlte Tranuk sich stets aus eisigen Augen beobachtet. Manchmal glaubte er, im Schneegestöber vor sich eine geisterhafte Gestalt mit ausgebreiteten Armen zu erkennen. Seit seiner Verwandlung fürchtete Tranuk sich vor keiner Auseinandersetzung mit dunklen Kreaturen oder uralten Zauberern, doch hatte er gehörig Respekt vor Geistern des Eises. Mit den verirrten Seelen im Schnee Gestorbener war nicht zu spaßen. Darum gebot er seiner gehörnten Steppenechse, rasch zurück ins sichere Cavern zurückzueilen, wo er seine wertvolle Fracht wieder in die Obhut der nun umso wachsameren Schildzwerge gab. So erwartete Tranuk die Rückkehr der restlichen Helden aus dem Land der drei Brüder. Erst, nachdem jene das geheimnisvolle Eiswesen unschädlich gemacht hätten, würde er sich wieder auf den Weg zur Höhle der Verwandlung machen. Dieses Abenteuer würde eines Tages als Legenden vom Herz aus Eis in die Chroniken der Bewahrer vom Baum der Lieder eingehen.\bigskip

Als die Helden wieder erwachten, befanden sich alle wieder zurück in ihrer eigenen Zeit. Es war ihnen, als hätten sie einen langen Schlaf mit wilden Träumen hinter sich gehabt. Und sie teilten nur noch unwirkliche, unzusammenhängende und uneinordbare Erinnerungen an die wilden Ereignisse der letzten Tage. 

Zunächst ziemlich verwirrt, kamen die Helden nach Augenzeugenberichten von aus der Ferne zugeschaut habender Andori zum Schluss, dass sie einen weiteren gefährlichen Gegner Andors erfolgreich besiegt hatten. Ein glühendes Wesen, welches durchs Land gehuscht war und seltsame Portale geöffnet hatte. Doch diese Portale waren nun geschlossen und die Gestalt vertrieben, auch wenn sie im finalen Kampf offenbar die Geister der Helden benebelt und ihre Erinnerungen verschleiert hatte. 

Die Helden in der Zukunft wurden zudem darüber informiert, dass sie die Gestalt gefangen und eingesperrt hatten. Leider hatte diese sich aus der Gefängniszelle unter der Rietburg wegteleportieren können, ehe sie die Helden hätte über ihre Motive informieren können. Doch viel wichtiger waren die Taten der restlichen Helden im Steppenland gewesen. Vögel zwitscherten. Die Eisschlafenden erwachten nach und nach: Gilda, Thoralds Trupp, Chada, Hallgard, Melkart und noch so viele mehr ... selbst der alte Troll neben dem alten Wehrturm! Die ewige Kälte lichtete sich. Der Frühling nahte. Die Legende hatte ein gutes Ende genommen. Alles war gut. 

\az{Jahr 63}

Orweyns glühende Gestalt stand in der dunklen Nacht neben der Taverne zum Trunkenen Troll. Sein langer Mantel flatterte im Herbstwind. Der rote Mond beschien seine finstere Taten. Er rieb sich seine schmutzigen Hände. Die zukünftigen Helden waren in ihre Zeit zurückgebracht und die vergangenen hatten immerhin ihre Erinnerung verloren. Und nun war sein wichtigstes Werk vollbracht. Unter ihm, tief vergraben im Schatten der andorischen Erde, ruhte ein unscheinbarer, bläulich schimmernder Quader, übersät von magischen Runen, der leise vor sich hin summte. 

Ein Winterstein. Orweyn hatte ihn hierher befördert, wie von seinem älteren Selbst aufgetragen.

Die Geister der andorischen Erde würden sich gegen den Eindringling in ihr Reich wehren. Doch selbst der große Skuvar würde Monate brauchen, um das magische Artefakt an die Oberfläche zu befördern, ohne selbst zu Eis zu erstarren. Und bis dahin würde sich die Ewige Kälte längst über Andor ausgebreitet haben. 

Der mächtige Zauberer vergoss eine Träne an all die Unschuldigen, die dieses Vorhaben kosten würde. Doch das war es wert, erinnerte er sich selbst. Der schlimme Tarok, die verwandelten drei Brüder, der wütende Enkel seines einstigen Freunds ... so viele lebensbedrohliche Gefahren lauerten in der Zukunft der bekannten Welt. Und auch wenn er es gerne anders hätte, sahen die Chancen nicht gut aus, dass die Helden sie alle überwinden konnten. Er musste realistisch sein. Drastische Maßnahmen waren gefordert. Sein Werk würde die Gefüge der Welt richten, wie es einst sein eigenes Opfer in Hadria getan hatte. Naja, sein eigenes Leben hatte er damals nicht geopfert, als er den Eisernen Turm hatte erbauen lassen. Das wäre ja auch unnötig gewesen. Nur die Leben der anderen Magier, denen er nicht vertrauen konnte, dass sie auch wirklich die Dunkle Magie hinter sich lassen würden. Und zumindest seine eigenen Erinnerungen an die Dunkle Magie würde er eines Tages opfern, wenn sein Ziel erreicht war und Frieden über dieser Welt lag. Nützliches kleines Tränklein, diese Tinktur gegen Seelenschmerz. 

Doch bis dahin gab es noch viel zu tun.

Orweyn würde nicht ruhen, bis die Dunkle Magie und ihre Bosheit endgültig aus dieser Realität getilgt waren. 

Vielleicht könnte er zurückkehren, wenn alles getan war, wenn Andor und das Land der drei Brüder unter meterdickem Schnee ruhten. Er könnte zur Knochengrube aufbrechen, Taroks dem Eisschlaf anheimgefallenen Körper finden und ihn durch gezielte Schwertstöße vernichten. Vielleicht auch noch die Hexe Reka und ähnlich gefährliche magische Zeitgenossen aus dem Wege befördern. Und danach die Wintersteine dorthin zurückzubringen, woher sie einst gekommen waren. Vielleicht würde der Schnee sich wieder lichten. Vielleicht würden einige wieder erwachen. Vielleicht. 

Er konnte hoffen.

Doch seine Taten waren nötig.

Niemand verstand das.

Niemand außer ihm.

Er war alleine in dieser kalten Welt.

Ein Retter, den keiner als Retter erkannte.

So wollte es das Schicksal von allen großen Rettern. 

Über dem Greis begann es zu schneien.

Die Ewige Kälte hatte begonnen.\bigskip

\az{63 und 64}

Doch waren damit wirklich alle losen Enden geschlossen? Wie war es eigentlich dazu gekommen, dass sich die Zeitportale gerade im passenden Moment wieder verschlossen hatten? Nachdem der jüngere Orweyn die zeitreisenden Helden in ihre eigene Zeit zurückgebracht hatte, aber bevor der ältere Orweyn in seiner Zelle aufgewacht war – und prompt mit einem Knall davongesprungen, um den Helden im Land der drei Brüder eine Standpauke über die Gefahren der Zukunft zu halten, mannigfaltig und grausam, wie er sie vorausgesehen hatte? Ein wenig früher, und manche Helden wären in der falschen Zeit gefangen gewesen. Ein bisschen später, und der ältere Orweyn hätte den Winterstein an sich gerissen. \bigskip

Nun, fürs Schließen der Zeitportale war ich verantwortlich. Nicht lange nach dem Endkampf tauchte ich neben dem Schlachtfeld auf. Dort sicherte ich die hadrische Tasche, die Orweyn in der Hitze des Gefechts verloren hatte und um die sich kein Held geschert hatte. Der Großteil des Inhalts war mir egal, den ließ ich liegen. Alte Bücher, wirre Worte auf brüchigem Papier, seltene Steine und magische Pülverchen aller Sorten. Orweyn würde sie später wieder an sich nehmen. Mir ging es nur um etwas: Die Sphäre aus Cavern, die das ganze Chaos hier ausgelöst hatte. Sie stotterte weiterhin wild vor sich hin und stieß blauen Zeitdampf aus. Mein Werk, mein armes Werk. Was hatten sie ihm nur angetan?! 

Mit geschickten Handbewegungen richtete ich verbogene verborgene Hebel.
 
Die Sphäre beruhigte sich. Die Zeitportale, welche wild über Andor gewirbelt hatten, schlossen sich. Die Zeit war nicht mehr entzweit.

Ich stellte sicher, dass 5 Gold und eine hilfreiche Nachricht beim Krallenfelsen lagen. Danach gab es nur noch etwas zu erledigen. Ich reiste ins herbstliche Cavern, stapfte verborgen durch enge Gänge und fand an der Stelle, die mein zukünftiges Selbst mir verraten hatte, einen zitternden Zwerg. Mart hieß er, und er war als einziger in der falschen Zeit verblieben. Argwöhnisch starrte er mich an. „Alles wird gut“, versprach ich, „Ich bin hier, um alles in Ordnung zu bringen.“ 

Ich zeigte Mart seine Sphäre, die ich repariert hatte und die ich nun in die geheime Schatzkammer zurückbringen würde. Und ich zeigte ihm meine zweite Sphäre, mit der ich danach in meine eigene Zeit zurückkehren würde. Dieselbe Sphäre, nur einige Jahrhunderte jünger als die andere. Mart machte große Augen. Ehrfürchtig fiel er auf ein Knie. Vermutlich hatte er mein Antlitz von den vielen Wandteppichen aus der Urzeit wiedererkannt. Leise flüsterte er: „Ihr kommt aus der Hochblüte des Zwergenreichs! Ihr seid ein Zeitreisender! Habt ihr die Zukunft gesehen? Werden wir die Ewige Kälte überleben?“ 

Ich lachte, zog ihn hoch und legte verschwörerisch einen Finger an meine Lippen. „Wer weiß, was die ferne Zukunft bringt. Schließlich kann ich so weit reisen, wie ich will, und ihr Horizont liegt weiterhin in unerreichbarer Ferne. Ich verirrte mich einst kurz in einer weit, weit entfernten Zeit, die jedem hierzulande unvorstellbar erschiene. Doch zumindest für die nächsten paar Jahrhunderte wird die Welt, so wie du sie kennst, florieren dürfen. Dafür sind wir Hüter der Zeit ja da.“ 

„Ihr Hüter der Zeit? Es gibt noch mehr?“

An einer Hand zählte ich die auf, die mir in den Sinn kamen: „Oh, es gibt so viele Hüter der Zeit! Ein tulgorischer, der sich darauf spezialisiert, den richtigen Leuten die richtigen Informationen zu geben, ehe er sie selbst vergisst. Ein Agren, der weiß, dass seine Höhlenmalereien Jahrhunderte später den Lauf der Dinge zu verändern vermögen, auch wenn seine Worte dann längst vergessen sind. Eine Drächin, die ihre eigene Urururgroßmutter ist. Ein grantiger Seher aus Danwar, der vor lauter Wald die einzelnen Bäume nicht mehr sieht. Ein zweiäugiger Zukunftswisser aus dem Süden, der vom ganzen metaphorischen Wald nur einen einzelnen Baum pflegt. Sie alle haben ihre Methoden. Oder hatten. Keine Ahnung, welche jetzt gerade überhaupt noch leben. Ich persönlich lasse sie alle ihre Spielchen spielen, davon verstehe ich nicht genug. Ich kümmere mich nur darum, dass meine Werke bis zu ihrer Zerstörung nicht zu lange in den falschen Händen bleiben. Das schulde ich der Welt dafür, die Sphäre geschaffen zu haben.“ 

Mart schienen tausend Fragen auf der Zunge zu liegen. Er entschied sich für eine bestimmte: „Wenn nicht Ihr ... gibt es denn irgendjemanden sonst, der den Überblick hat?“ „Wer kann das schon wissen? Mutter Natur, die Schicksalssporne, Mächte jenseits unserer Vorstellung? Falls es sie gibt, haben sie sich noch nicht bei mir gemeldet. Aber falls sie wirklich einen solchen Durchblick hätten, müssten sie sich kaum direkt bei mir melden. Wir können nur hoffen, dass sie unsere Werte teilen.“ Ich setzte mein liebstes bevormundendes Lächeln auf. „Die wichtigsten Hüter der Zeitlinie sind natürlich ihre Bewohner, die sich tagein, tagaus mit Elan für das Gemeinwohl einsetzen: Bauern und Bewahrer. Heiler und Helden. Fackelträger aus den Tiefminen und Fürsten von Cavern gleichermaßen. Du. Die Deinen. Ihr Werk können wir tagein, tagaus spüren und bejubeln. Doch hin und wieder tauchen Gefahren auf, deren Ausmaß diese Hüter kaum einschätzen können. Jahrhunderte lang geplante Pläne von finsteren Sehern und erbarmungslosen Mächten. Perfide Prophezeiungen und stibitzte Sphären. Dann gibt es oft lose Enden zu schließen. Irgendjemand muss schließlich dafür sorgen, dass die Zeitlinie konsistent und gut ist. Ich habe mich entschieden, meinen Teil zu tun.“ 

Mart schien sich nicht sicher zu sein, ob er erleichtert oder angegruselt sein sollte.

Behutsam führte ich ihn zurück zu seinem Heim in seiner eigenen Zeit, in jener der Schnee der Ewigen Kälte langsam zu schmelzen begonnen hatte. Mart lud mich auf eine Tasse heißen Cocoa hinein, doch ich lehnte höflich ab. Ich wollte nicht zu viel Zeit mit Plaudereien verbringen, erst recht nicht so weit abseits meiner eigenen Zeit. So blieb ich zurück, als Mart die Tür zu seiner Höhle öffnete. Ich vernahm noch, wie Marts Lebenspartner Bort voller Freude und Sorge zugleich aufjubelte, als er Mart durch die Pforte treten sah: „Bei Boords Bart, ich hatte mir schon solche Sorgen gemacht! Glaube mir, ich war so kurz davor, mit deiner Schwester Ragga aufzubrechen, diesen glühenden Greis eigenhändig aufzuspüren und in die Eisenketten meines Vetters zu legen! Wenn Drak mich nicht beruhigt hätten ... Oh, welch Horror du erlebt haben musst! Wie geht es dir?“ 

Leise flüsterte Mart eine beruhigende Entgegnung.

Und ich spionierte den beiden nicht länger nach, sondern öffnete ein Portal viele Jahrhunderte in die Zukunft, wo der nächste Auftrag meines zukünftigen Selbsts – und eine Bewahrerin namens Josella – bereits auf mich warteten. 

